\Needspace{8\baselineskip}\originalchapter{输运、鞘层与反磁响应}\label{ch:transport}
\originalsection{随机游走和扩散}
碰撞使速度相关性在时间 $\tau\sim1/\nu$ 衰减. 沿磁场典型步长 $v_{\rm th}\tau$, 扩散系数 $D_\parallel\sim v_{\rm th}^2/\nu$. 强磁化横向每次碰撞改变导引中心约一个 $r_L$, 因而 $D_\perp\sim\nu r_L^2$. 两者在 $|\Omega|\gg\nu$ 时相差约 $(\Omega/\nu)^2$. 这是经典输运标度, 湍流和环形新经典输运可产生更大横向损失.

无 B 的摩擦动力学 $m\nu\mathbf u=q\mathbf E-\nabla p/n$ 给迁移率和扩散关系 $D=T/(m\nu)$. 电子本来扩散较快, 会使边界电势建立并限制电子、加快离子. 准中性、等通量条件下产生双极扩散, 不能让每种粒子永远独立按自己的扩散系数流失.

\originalsection{单向热通量}
对各向同性速度分布$f(v)$, 入射平面的单向通量为$\Gamma=\int_{v_x>0}v_xf\dd^3v$. 取以法线为极轴, 角积分和结果为
\begin{align}
 2\pi\int_0^{\pi/2}\cos\theta\sin\theta\dd\theta&=\pi,\\
 \Gamma&=\frac n4\langle v\rangle.
\end{align}
Maxwell 分布 $\langle v\rangle=\sqrt{8T/(\pi m)}$, 所以 $\Gamma=n\sqrt{T/(2\pi m)}$. 沿单坐标速度均值零并不意味着平面上单侧无入射通量, 因为另一侧反向通量仅在净流中抵消.

\originalsection{鞘层与空间电荷限制}
壁面附近准中性失效. 热电子流失快使浮动壁通常变负, 减少电子通量. 冷离子穿入稳定鞘需达到声速量级, 简单模型 Bohm 条件 $u_i^2\ge T_e/m_i$. 它依 Maxwell电子、冷离子和单调鞘层假设, 不把某个常数声速作为所有多温、多离子模型的通用公式.

无离子的平板真空二极管, 从零速阴极到阳极势 $\phi_A$, 能量守恒 $mv^2/2=e\phi$, 电流密度幅值 J满足 $n=J/(ev)$, Poisson 方程 $\phi''=J\sqrt{m/(2e)}\phi^{-1/2}/\epsilon_0$. 两边乘 $\phi'$ 积分, 最大可通过电流对应阴极场零. 再积分给 Child--Langmuir 定律
\begin{equation}\label{eq:child}
 J_{\max}=\frac{4\epsilon_0}9\sqrt{\frac{2e}m}\frac{\phi_A^{3/2}}{x_A^2}.
\end{equation}
零速边界使密度形式上在阴极发散, 是理想模型; 非零注入速度改变边界和系数.

\originalsection{导引中心电流和磁化电流}
磁矩运动形成微观环流, 宏观磁化 $\mathbf M=-p_\perp\mathbf B/B^2$. 导引中心漂移算出的电流 $\mathbf J_{
\rm gc}$ 不总等于全粒子流电流, 还需 $\nabla\times\mathbf M$. 将两者相加才可同 Ampère 方程中的全电流比较. 这解释反磁漂移为何可在无共同物质输运的情况下产生压强梯度电流.

各向同性流体的反磁速度 $\mathbf v_{*s}=\mathbf B\times\nabla p_s/(q_sn_sB^2)$, 总反磁电流为 $\mathbf J_*=\mathbf B\times\nabla p/B^2$. 它是局部速度矩的表达, 不应一律当作整团粒子的导引中心质量流. 由场曲率产生的粒子漂移和磁化修正同样要分开计数, 后面 z-pinch 作业具体检验这一抵消.
